\section{辅政、淮南与再统一}
\label{sec:three-kingdoms-regencies-reunification}

景初三年，曹叡去世，幼帝曹芳即位，曹爽与司马懿受遗诏辅政。三国中期的关键，不再只是三位皇帝如何宣称自己的合法性，而是幼主朝廷如何让禁军、诏令和大臣任用继续运转。遗诏的全文、近侍的作用和两位辅政者最初的权力边界都有争议；可以确知的是，魏的皇帝中心此后必须经由一套容易被争夺的辅政结构才能发挥作用。

\subsection{洛阳的辅政与江东的继承}

正始年间，魏廷在洛阳太学立三体石经，以古文、篆、隶校写经典。它同二三八年倭女王卑弥呼遣使入魏一样，显示洛阳仍能以学校、册封和文字向内外表达朝廷的权威；石经的篇目、完成阶段以及册封所到达的实际关系范围，却都不能由一条记载推成无缝的文化或政治统属。曹爽二四四年攻蜀至兴势而退，也再次说明，即使中枢能调集军队，山地、守险和粮道仍会限制进攻。

孙吴面对的不是同一套辅政人事，而是继承人彼此形成的宾客网络。孙登死后，孙和被立为太子，孙霸封鲁王；二五〇年孙权废孙和、赐死孙霸，改立幼子孙亮。废立和死亡的结果可靠，参与者的罪责和每一次宫廷言辞却多出自后来的传记。把二宫之争写成单纯的家事，会遮蔽它如何清洗官僚并削弱成熟的继承安排。二五二年孙权死，诸葛恪受遗诏辅政，江东也进入由大臣的军政表现来检验幼帝朝廷的阶段。

\subsection{高平陵以后，皇帝仍在而中枢已变}

二四九年司马懿在高平陵政变中控制洛阳，曹爽失势并随后被杀。政变日的军队部署、诏令真伪和司马懿的承诺是否构成欺罔，不能由《三国志》与《资治通鉴》的后出叙事完全裁定；其制度结果却清楚可见：司马氏取得了魏廷的主要军政与人事，曹魏皇帝仍在位，却不再能以原来的方式制约辅政者。

这种格局在淮南反复接受考验。二五一年王凌谋立楚王曹彪失败；二五五年毌丘俭、文钦起兵；二五七年至二五八年诸葛诞据寿春反司马昭并向孙吴求援，最终城陷。三次行动的动机、响应者和兵额各不相同，不能合并成一场同质的“反司马氏运动”。它们共同显示，淮河一线连接魏、吴与北方军镇，任何一支掌兵集团都可能以皇帝、宗室或盟国的名义挑战洛阳。司马昭的胜利并非消除了这个结构，而是把掌握中军和镇压外镇的能力更集中地放到司马氏手中。

孙吴的权臣政治则继续侵蚀皇帝中心。诸葛恪围合肥失利后被孙峻所杀，孙峻死后孙綝接掌禁军；二五八年孙綝废孙亮，另立孙休。吴的水陆防务仍能在濡须等方向抵挡魏军，但宫廷连番废立使将领、宗室和官僚难以获得稳定的调度中心。不能以一位皇帝或权臣的品行解释这一切；边防、幼主、禁军和宾客集团彼此牵动，才是江东政治越来越难自我修复的条件。

\subsection{从灭蜀到灭吴}

二六三年魏多路伐蜀，邓艾越阴平抵成都平原，刘禅出降；翌年钟会在成都图谋自立而败，表明吞并一个政权并不等于立即取得其地方军政网络。阴平行军的筹划、成都的守备与谯周等人的作用，各有难以化解的史料争议；蜀汉终结、降诏改变姜维等前线将领的行动空间，则是更可靠的事实。灭蜀后，司马昭加九锡、进晋公又进晋王，封国与魏廷之间的界线更加空洞。

二六五年曹奂禅位于司马炎，西晋建立。禅让诏册如何撰写、群臣如何被动员、曹奂是否有实质选择，均应同既成的军政控制分开处理。西晋颁行《泰始律》，又在荆州、淮南经营屯戍与水军；羊祜、杜预和王濬等人的部署，使二七九年的伐吴能够沿江、荆州与淮南多路推进。王濬舰队抵建业、孙皓出降的结果结束了三国并立，但功绩归属、降仪和战后安置仍有争议。统一首先是一个军事和政权接收的结果，不是江南社会、赋役和地方网络立刻变成同一尺度的证明。

\subsection{三国末年的材料边界}

《三国志》是魏、蜀、吴宫廷与战争年序的主干，陈寿以魏为轴的编排影响称谓和叙事篇幅；裴松之注保留丰富异说，却不能与正文视作同一层级。《晋书》可校看魏晋禅代、泰始律与伐吴，仍带有统一之后的回望。走马楼吴简、三体石经残石和区域考古让基层行政、文本与空间不只留在帝纪中，但它们各自的地域、断代和保存条件限制很强。本节只用这些材料说明辅政、外镇、道路与继承如何共同瓦解三国的均势，不把兵数、心迹或疆域写成确定无疑的结论。
